\documentclass[11pt,a4paper]{article}
\usepackage{kagan}

\begin{document}
\phaseheader{15}{Hardening, Polish and Release: Kagan 1.0}{Make it survive hostile input, measure it, document it, ship it, and choose where to go next}{96}{100}

\ataglance{3--4 weeks for the core, then open-ended}{4}{Phases 0--14 complete: Kagan boots to a shell with pipes, files and processes}{A
system call fuzzer that runs for an hour without a kernel panic, an audited user/kernel boundary, kernel stack
protection and read-only kernel code, users, file and device permissions with sessions, a threat model with its
residual risks, a benchmark suite with one measured optimisation, an architecture document, a reproducible
\code{make release} that produces a bootable ISO with published hashes, the signed \code{v1.0} tag, and a
written design review. After that, optional extension tracks, and Stage V: the RISC-V and ARM
ports of Phases 16--18.}

\section{Scope and motivation}

A system isn't finished when its features work on the inputs it was tested with. It is finished when it survives inputs
nobody intended, when its performance is measured and explained, when someone else can build and understand it, and when
every design choice in it is documented and justified. That is the difference between a demo and version 1.0, and this phase is about
closing that gap.

The release is also the point for a design review. \MOS chapter 12, written by an author who designed several operating
systems, covers exactly the decisions made across Phases 0--14: interfaces, structure, mechanism versus policy,
performance, and managing a large software project. The phase ends with a written comparison of Kagan against
MINIX 3 and Linux, a list of optional extension tracks, and a look ahead to Stage V, which ports Kagan to RISC-V
(including the project's own hardware core) and to ARM.

\section{Background}

\subsection{The kernel is the trusted computing base}

\MOS\,\S9.2.2 defines the \textbf{trusted computing base} (TCB): the hardware and software that must be correct for
the system's security to hold. In Kagan that is the entire kernel. Every system call is an entry point into the TCB
controlled by a possibly hostile program, which is why Phase 10's rule (never trust a user pointer) was only the
start. \MOS\,\S9.7 catalogues how kernels get broken, and several items apply directly to Kagan code:
\textbf{buffer overflows} when a length from the user isn't checked; \textbf{integer overflows} when
\code{offset + length} wraps around; \textbf{null-pointer dereferences} in the kernel; and \textbf{time-of-check to
time-of-use} (TOCTOU) races, where the kernel checks a value in user memory and then reads it again after the user has
changed it. The cure for the last one is to copy first and check the copy.

\subsection{Protection: who may do what}

Protection needs a model of \emph{subjects} (users, processes), \emph{objects} (files, devices) and \emph{rights}
(read, write, execute). \MOS\,\S9.3 and \OSC chapter 17 describe the two ways to store it: \textbf{access control
lists} attached to objects, and \textbf{capabilities} held by subjects. UNIX uses a compact ACL, nine permission bits
for owner, group and others in each inode, plus the rule that user ID 0 (root) may do anything. KaganFS already stores
the mode, user ID and group ID in every inode (Phase 13), so what is left is to give processes an identity and check it.

\subsection{Measuring before optimising}

\MOS\,\S12.4 gives the rules for performance work: find out what is actually slow, optimise the common case, use
caching and hints, and exploit locality. None of that works without measurements. The x86 \code{rdtsc} instruction
reads a cycle counter, which is precise enough to time a single system call. The classic operating-system
microbenchmarks (in the style of \emph{lmbench}) measure the null system call, a context switch between two processes
passing a token through pipes, process creation, and file-system throughput.

\subsection{Design, looking back}

\MOS\,\S12.2--12.3 lists principles to measure Kagan against: \textbf{simplicity} and
\textbf{completeness} of the interface, \textbf{orthogonality} (independent concepts that combine freely, like
\code{fork} and \code{exec}), \textbf{separation of mechanism and policy} (the Kagan scheduler), static versus dynamic
structures (MINIX's fixed tables versus the Kagan heap), and top-down versus bottom-up implementation (Kagan is
built bottom-up).
\S12.5 adds Brooks's observations about software projects, including the \emph{second-system effect}: the temptation to
put everything left out of version 1 into version 2.

\section{References}

\begin{readingplan}
\must & \MOS & Ch.~12 \emph{Operating System Design}: nature of the problem, interface design, implementation, performance, project management, trends & 981--1030 & The design principles behind the Task~15.8 review. \\
\must & \MOS & \S9.1--9.3 \emph{The Security Environment}, \emph{Operating Systems Security}, \emph{Controlling Access to Resources} & 595--611 & Threats, the TCB, protection domains, ACLs and capabilities. \\
\must & \MOS & \S9.7 \emph{Exploiting Software} (buffer overflows, format strings, dangling and null pointers, integer overflows, TOCTOU) & 639--657 & The bug classes the boundary audit looks for. \\
\should & \MOS & \S9.10 \emph{Defenses} (especially code signing and jailing) & 684--703 & What operating systems add on top of correct code. \\
\should & \OSC & \S16.1--16.2 \emph{The Security Problem}, \emph{Program Threats}; \S16.5--16.6 \emph{User Authentication}, \emph{Implementing Security Defenses} & 621--634, 648--662 & A second view, with passwords and defences in depth. \\
\should & \OSC & Ch.~17 \emph{Protection}, \S17.1--17.6 (principles, rings, domains, access matrix) & 667--682 & The formal model behind UNIX permissions. \\
\should & \OSDI & \S5.4--5.5 \emph{Security}, \emph{Protection Mechanisms} & 526--548 & Compact coverage of authentication, ACLs, capabilities and covert channels. \\
\should & \OSDI & \S1.5 \emph{Operating System Structure}; \S2.5.1 \emph{The Internal Structure of MINIX 3} & 42--51, 112--116 & Monolithic versus microkernel, for the design review. \\
\should & \MOS & \S10.1--10.8 \emph{Case Study: UNIX, Linux, and Android} (skim) & 713--848 & The Linux column of the design review. \\
\could & \OSC & Ch.~20 \emph{The Linux System} & 775--819 & Another view of Linux for the comparison. \\
\should & \LKD & Ch.~18, \emph{Debugging} & 363--378 & Oopses, assertions and debugging techniques for the hardening work. \\
\could & \LKP & Ch.~5, \emph{Kernel modules and security -- an overview}; Ch.~7, \emph{Randomizing the memory layout -- KASLR} & 256--261, 362--367 & Defences a production kernel adds. \\
\should & \SPEC & NIST SP~800-218, \emph{Secure Software Development Framework} (\url{https://csrc.nist.gov/pubs/sp/800/218/final}), practices PS and PW & -- & The yardstick for Task~15.7's release integrity and Task~15.2's review. \\
\should & \SPEC & Linux kernel documentation, \emph{Kernel Self-Protection} (\url{https://docs.kernel.org/security/self-protection.html}) & -- & The defences of Task~15.3, and the ones Kagan leaves out. \\
\end{readingplan}

\begin{minixbox}
MINIX 3's central design goal is reliability, and \OSDI\,\S2.5.1 explains how its structure serves it. Drivers run
as separate user processes with only the privileges they need, so a buggy driver can't overwrite the kernel. A
special process, the \textbf{reincarnation server}, watches the drivers and servers and restarts any that crash
(\OSDI pp.~114 and 599). In Kagan, every driver bug is a kernel bug. This contrast belongs in the design review of Task~15.8,
and Track E below applies the MINIX approach to one Kagan driver.

\par\smallskip\textbf{Linux, for comparison.} \LKD\ \LKP Linux's hardening goes further: KASLR, signed modules, lockdep and memory-error detectors such as KASAN, plus the \code{BUG_ON}/\code{WARN_ON} assertions and oops reports that \LKD's chapter on debugging describes.
\end{minixbox}

\section{Implementation}

\task{Fuzz the system call interface}
\code{user/bin/fuzz} has made random system calls since Phase 11, with a seed you can print and replay, and Phase 14
taught it pipes and the terminal. This is its final pass. Check that the argument mix covers every class: small
integers, huge integers, negative numbers, pointers to valid buffers, pointers into the kernel (\code{0xC0000000} and
up), pointers into the program's own read-only text, unmapped addresses, addresses one byte before the end of a page,
descriptors that exist and ones that don't, and random strings as paths. Add each call of this phase as you write it.
Run several copies at once together with a stress script that forks, execs, pipes and writes files in parallel, first
as root and, once Task~15.4 exists, as an ordinary user (UID 1000). Every kernel panic is a bug to fix; every process
killed by a signal is fine. If you have a spare machine, also let the campaign soak for 24 hours locally; hosted CI
jobs are time-limited, so the long run stays out of CI.
\verify{Four fuzzers and the stress script run for one hour with zero kernel panics and no growth in memory use
after they are stopped, both as root and as UID 1000.}

\task{Audit the boundary}
The boundary table \code{docs/security/syscalls.md} has grown since Phase 10, one row per system call with each
argument and how it is validated. Make the final pass: every call, including those of this phase, has a complete row,
the code still does what the row says, every user pointer goes through the uaccess routines, and \code{cppcheck} runs
clean. Then look specifically for these bugs, because the fuzzer is unlikely to find all of them.

\begin{center}
\small\renewcommand{\arraystretch}{1.25}
\begin{tabularx}{\linewidth}{@{}l X@{}}
\toprule
\tkth{Bug class} & \tkth{Where to look in Kagan} \\
\midrule
missing pointer check & any use of a user pointer outside \code{copy_from_user}/\code{copy_to_user}/\code{strncpy_from_user} \\
integer overflow & \code{offset + len} in \code{read}/\code{write}/\code{lseek}, \code{brk} arguments, ELF segment sizes, \code{argv} totals \\
TOCTOU & paths or structures read twice from user memory; copy once into the kernel, then validate the copy \\
unchecked index & descriptor numbers, signal numbers, system call numbers, directory entry indices \\
information leak & kernel stack or heap bytes copied to the user in padding of \code{struct stat} or \code{dirent}; zero structures before filling them \\
resource exhaustion & unlimited processes, descriptors, pipe buffers or path lengths; every resource needs a limit \\
\bottomrule
\end{tabularx}
\end{center}

\task{Harden the kernel itself}
The kernel has been built with \code{-fstack-protector-strong} since Phase 10, and its \code{.text} and \code{.rodata}
have been read-only since Phase 6. Check both once more on the release build: every kernel object file gets the
protector except the one function that sets \code{__stack_chk_guard}, which is compiled with
\code{-fno-stack-protector} and never returns into a frame entered before the guard changed (that return would report
``stack smashing'' on a correct stack). The kernel image is still mapped with 4\,KiB pages, CR0.WP is still set, and
the page at address 0 is never mapped in any process.

The guard is only as good as its seed, and \code{rdtsc} alone is predictable. Give the kernel one entropy function:
\code{rdrand} when CPUID reports it (QEMU's \code{-cpu max} does), otherwise \code{rdtsc} mixed with the Multiboot
memory map. Write in the threat model that the fallback is weak and acceptable only for a teaching system.

\begin{lst}{c}{kernel/arch/i386/random.c}
#include <cpuid.h>

uint32_t arch_random32(void)
{
    unsigned a, b, c, d;
    if (__get_cpuid(1, &a, &b, &c, &d) && (c & bit_RDRND)) {
        for (int i = 0; i < 10; i++) {          /* rdrand may fail briefly; retry */
            uint32_t v;
            unsigned char ok;
            __asm__ volatile("rdrand %0; setc %1" : "=r"(v), "=qm"(ok) :: "cc");
            if (ok)
                return v;
        }
    }
    return weak_seed();                         /* rdtsc ^ memory map hash: predictable */
}
\end{lst}

Two CPU features cost one bit each. When CPUID leaf 7 reports \textbf{SMEP} (EBX bit 7), set CR4 bit 20: the CPU then
refuses to execute a user page in ring 0, so a corrupted kernel function pointer can't jump into code a program wrote.
\textbf{SMAP} (EBX bit 20, CR4 bit 21) is optional. It also blocks ring-0 data access to user pages, so the uaccess
routines must wrap their copy in \code{stac} and \code{clac}; with a single uaccess choke point that is a two-line
change.

Security denials should leave a trace. Log a failed \code{setuid}, a permission denied on a device node and a process
killed by a fault with \code{pr_warn} and the prefix \code{security:}, at that one fixed level, so
one search of \code{log_buf} for \code{security:} answers ``what was refused?''. Finally, make crash reports readable: generate a symbol
table from \code{i686-elf-nm} at build time, load it as a boot module, and print function names in the backtrace of
every panic.
\verify{A deliberate overflow of a local buffer in a test system call panics with ``stack smashing detected'' and a
symbolised backtrace; a write to a function's code faults. Under \code{-cpu max} the boot log reports \code{rdrand} and
SMEP, and a test system call that jumps to a user address faults instead of running it.}

\task{Users and permissions}
Give each task a user ID and a group ID, inherited through \code{fork}. Add \code{getuid} (24), \code{setuid} (23),
\code{getgid} and \code{setgid}; only root may change identity. Fill in \code{vfs_permission} (Task~13.1), which every VFS operation already calls,
and check permissions there against the inode's mode,
owner and group on \code{open}, \code{execve}, \code{create}, \code{mkdir} and \code{unlink}, with root bypassing the
checks. The same checks cover \code{rename}, \code{link}, \code{chmod}, \code{chown} and \code{truncate}, including
\code{open} with \code{O_TRUNC}; only the owner or root may \code{chmod}, and only root may \code{chown}.

Files are not the only objects. A user who can open \code{/dev/hda} can rewrite every inode, \code{/etc/passwd}
included, and walk past all the checks above. Give each \code{devfs} node an owner, a group and a mode, and check them
through the same VFS path: \code{/dev/hd*} is \code{root:root 0600}, \code{/dev/tty} is \code{0620}, and
\code{/dev/null} and \code{/dev/zero} are \code{0666}.

The terminal needs an owner too. As Phase 14 left it, \code{ioctl} 54 lets any process make any group the foreground,
so a background job left behind by one user could take the console and read the next user's \code{login} input. Give
each process a session ID, inherited through \code{fork}, and add \code{setsid}, which makes the caller the leader of
a new session. The first terminal a session leader opens becomes the session's \textbf{controlling terminal}.
\code{ioctl} 54 then accepts only a process group in the caller's own session and returns \code{-EPERM} otherwise.
When the session leader exits, the kernel sends \code{SIGHUP} to every process in the session and revokes their
descriptors on the terminal, so later reads and writes on them fail.

Add \code{/etc/passwd} with lines of the form \code{name:uid:gid:home:shell}, and a \code{login} program that
\code{init} runs on the console, which looks up the user, starts a new session, changes to their identity and home
directory, and runs their shell. \code{login} opens \code{/etc/passwd} and every other file of its own with
\code{O_CLOEXEC}, so no descriptor opened as root survives into the user's shell. It drops privileges in the safe
order, group first, because after \code{setuid} it may no longer change its group, and it stops if either call fails.

\begin{lst}{c}{user/bin/login.c (dropping privileges)}
if (setgid(pw->gid) < 0 || setuid(pw->uid) < 0) {
    fprintf(stderr, "login: cannot change identity\n");
    exit(1);                    /* never run the shell as root by accident */
}
\end{lst}

Then decide about passwords and write the decision down. One honest option is no passwords at all: \code{login} only
asks for a name, and the threat model states that whoever has the console is trusted, which is true of a QEMU window.
The other is real passwords, and then never write your own hash. Port the reference code of a deliberately slow,
salted key-derivation function (bcrypt, scrypt or Argon2) into user space, compute the hashes on the host when you
build the image, and store them in \code{/etc/shadow}, owned by root with mode \code{0600}, because \code{/etc/passwd}
must stay readable by every user. A plain salted SHA-2 hash is too fast and is brute-forced in hours.
\verify{An ordinary user can't read a file with mode \code{0600} owned by root, can't write to \code{/bin}, and can't
kill root's processes; root can do all three. As UID 1000, every evil program still behaves as listed, \code{open} of
\code{/dev/hda} returns \code{-EACCES}, and a background job's \code{ioctl} 54 on another session's terminal returns
\code{-EPERM}. After logout, a job left running in the old session has received \code{SIGHUP}.}

\task{Benchmarks, and one real optimisation}
Write \code{user/bin/bench}, which reports the median of many runs for each measurement below, and put the results in
a table in your documentation. Then pick the worst number, find the cause with measurements, fix it, and record the
before and after values.

\begin{center}
\small\renewcommand{\arraystretch}{1.25}
\begin{tabularx}{\linewidth}{@{}l X@{}}
\toprule
\tkth{Benchmark} & \tkth{How} \\
\midrule
null system call & time 100\,000 calls to \code{getpid} with \code{rdtsc} \\
context switch & two processes pass a byte back and forth through two pipes; divide by the number of switches \\
process creation & \code{fork} + \code{exit} + \code{wait}; and \code{fork} + \code{exec} + \code{wait} of a tiny program \\
file throughput & write and read a 16\,MiB file on KaganFS, with a cold and a warm buffer cache \\
boot time & \code{rdtsc} in the loader and at the first shell prompt, printed on the serial port \\
\bottomrule
\end{tabularx}
\end{center}

Typical optimisations at this point are global pages for kernel mappings (no TLB flush of kernel entries on a CR3
switch), multi-sector disk transfers and read-ahead, a hash table in the buffer cache, and avoiding the double copy
through a kernel buffer in pipes. Run the same benchmarks in a Linux system for perspective; the point is to
understand the differences, not to win.

\task{Write the documentation}
Write an architecture document (a LaTeX file next to this plan, or Markdown): an overview diagram like the one in the
roadmap; each subsystem in a page or two; the virtual and physical memory maps; the system call table with numbers,
arguments and errors; the lock order; the KaganFS on-disk format; how to build, run, test and debug; and known
limitations. Add a \code{README} that gets a newcomer from a clean clone to a running shell, and a
license (MIT).

Bring \code{docs/security/threat-model.md} to version 1. Next to the assets, adversaries and trust boundaries it
already lists, add a table of residual risks: each threat, the control that addresses it, and what is left, such as
the weak entropy fallback, the missing W\textasciicircum{}X for user data on i386 without PAE, or console access trusted when there
are no passwords. Close it with an explicit out-of-scope list: physical attackers and the boot chain (BIOS boot has no
root of trust), network attackers (there is no network stack), malicious hardware and DMA, and speculative-execution
side channels such as Spectre and Meltdown. Writing them down is what separates a decision from an oversight. Add
\code{SECURITY.md} at the top of the repository: which version is supported, how to report a bug privately, and a link
to the threat model.

\task{Release v1.0}
Add \code{make release}: build the kernel, the user programs, the root archive and a KaganFS image populated with
\code{/bin}, \code{/etc} and a home directory, and a GRUB ISO whose \code{grub.cfg} loads the kernel and its modules.
Check that the release builds from a clean clone (\code{git clone} into a temporary directory, then \code{make release}),
boot the ISO in QEMU, and tag the commit \code{v1.0}. If you want to try real hardware, an older PC with legacy BIOS (or
UEFI with CSM) can boot the ISO from a USB stick.

\begin{warnbox}[Writing a USB stick]
\code{sudo dd if=build/kagan.iso of=/dev/sdX bs=4M} erases the target device completely. Run \code{lsblk} before and
after plugging the stick in, make sure \code{sdX} is the stick and not your laptop's disk, and never guess.
\end{warnbox}

``Reproducible'' is a claim you can test. Fix the time the build embeds with \code{SOURCE_DATE_EPOCH}, keep the
clone's path out of debug information with \code{-ffile-prefix-map}, build twice from two clean clones, and compare the
hashes. \code{kernel.elf} and the user binaries must be identical. \code{grub-mkrescue} may embed timestamps in the ISO;
if two ISOs differ, say so in the release notes instead of hiding it. Publish a \code{SHA256SUMS} file with the ISO, the
kernel and the KaganFS image, and sign the tag so anyone can check that it came from you.

\begin{lst}{sh}{Two clean builds, hashes and a signed tag}
export SOURCE_DATE_EPOCH=$(git log -1 --format=%ct)
git clone -q . ../rel-a && make -C ../rel-a release
git clone -q . ../rel-b && make -C ../rel-b release
sha256sum ../rel-a/build/kernel.elf ../rel-b/build/kernel.elf    # the two hashes must match
(cd ../rel-a/build && sha256sum kagan.iso kernel.elf kaganfs.img > SHA256SUMS)
git tag -s v1.0 -m "Kagan 1.0"                                  # needs a GPG key
\end{lst}
\verify{A fresh clone builds with one command and boots from ISO to the login prompt. \code{make test} passes on the
tagged commit. Two clean builds give identical \code{kernel.elf} hashes, \code{SHA256SUMS} matches the published files,
and \code{git tag -v v1.0} reports a good signature.}

\task{Write the design review}
Write a design review and architecture retrospective of three to five pages. For each subsystem (boot, memory, processes,
scheduling, synchronisation, system calls, file system, drivers), state what Kagan does, how MINIX 3 does it
(\OSDI) and how Linux does it (\MOS chapter 10), and what should change with hindsight. Use \MOS chapter 12's
principles as the yardstick, and end with the decision that held up best and the one to reverse first. Record the
review in \code{docs/} next to the architecture document, and walk through it with at least one reviewer who did not
write the code.

\section{Extensions after v1.0}

Kagan 1.0 is a complete small system. The planned continuation is Stage V (Phases 16--18), which ports it to two
more architectures with the same structure of tasks, checks and exit criteria. The tracks below are open-ended
instead. They extend Kagan in different directions, vary widely in size, and can be combined with each other and
with Stage V.

\begin{center}
\small\renewcommand{\arraystretch}{1.35}
\begin{longtable}{@{}>{\raggedright\arraybackslash}p{0.19\linewidth} >{\raggedright\arraybackslash}p{0.46\linewidth} >{\raggedright\arraybackslash}p{0.29\linewidth}@{}}
\toprule
\tkth{Track} & \tkth{Scope} & \tkth{References} \\
\midrule
\endfirsthead
\toprule
\tkth{Track} & \tkth{Scope} & \tkth{References} \\
\midrule
\endhead
\bottomrule
\endlastfoot
A.\ Graphics and Turkish text & Request a linear framebuffer through the Multiboot video fields, draw text with a PSF font that covers the whole Turkish alphabet (ğ, ş, ı, İ, Ğ, Ş), add a PS/2 mouse (IRQ 12), then a simple window system & \MOS\,\S5.6 (user interfaces); OSDev Wiki (VESA, PSF, mouse) \\
B.\ Multiprocessor (SMP) & Parse the ACPI MADT, program the local APIC and I/O APIC, start the other CPUs (INIT--SIPI--SIPI), per-CPU data and run queues, TLB shootdown; the Phase 9 spinlocks see real contention & \MOS\,\S8.1 (520--544); \OSTEP Ch.~10 (97--108); \OSC\,\S5.5 (220--227) \\
C.\ Networking & PCI enumeration, an e1000 or RTL8139 driver in QEMU, Ethernet, ARP, IPv4, ICMP (\code{ping}), UDP, then a minimal TCP and a socket API & \OSC Ch.~19 (733--768); \MOS\,\S8.3 (566--587); OSDev Wiki and RFCs \\
D.\ 64-bit port & Long mode, four-level paging, Multiboot2, \code{syscall}/\code{sysret}, the x86-64 calling convention. Much easier after Stage V: Phase 16 makes the generic code 64-bit clean, and Phase 18 builds a 64-bit MMU on AArch64 & \OSC\,\S9.6 (379--383); Intel SDM Vol.~3A \\
E.\ Microkernel experiment & Move the ATA driver (or KaganFS) into a user-space server that talks to the kernel with synchronous messages, as MINIX does, and measure the cost & \OSDI\,\S2.5, \S3.4 (112--125, 252--261); \MOS\,\S1.7.3 (65--68) \\
F.\ Virtualisation & Examine what QEMU and KVM provide, run Kagan with \code{-enable-kvm}, and compare the benchmarks & \OSC Ch.~18 (701--729); \MOS Ch.~7 (471--516) \\
G.\ Loadable kernel modules & Load relocatable ELF objects (\code{.ko}-style) into a running kernel: resolve their symbols against an exported symbol table, apply relocations, call \code{init} and \code{exit}, and move one driver (the RAM disk) into a module, with \code{insmod}, \code{rmmod} and \code{lsmod} commands & \LKP Ch.~4--5 (143--265); \LKD Ch.~17 (337--362) \\
\end{longtable}
\end{center}

Stage V deserves a note. Porting exposes every place where x86 leaked into supposedly generic code, and fixing those
leaks improves the whole kernel, including the x86 version. Stage V also runs Kagan on the project's own RISC-V
processor, so the kernel and the hardware are validated against each other. Stage V targets Kagan 2.0; this phase
closes with Kagan 1.0 complete.

\section{Exit criteria}

\begin{milestonebox}[Kagan 1.0]
\begin{checklist}
  \item One hour of parallel fuzzing and stress testing, as root and as UID 1000, without a kernel panic or a memory
        leak.
  \item Boundary audit table complete; stack protector, read-only kernel code and symbolised backtraces in place.
  \item Entropy source, SMEP when available, and \code{security:} log lines for denials in place.
  \item Users, \code{login} and file permissions enforced; device nodes owned, sessions and terminal ownership
        enforced; the password decision written down.
  \item Benchmark table written, with one optimisation measured before and after.
  \item Architecture document, README and system call reference written; threat model v1 lists every residual risk;
        \code{SECURITY.md} in place.
  \item \code{make release} works from a clean clone; the ISO boots; two clean builds give identical
        \code{kernel.elf} hashes; \code{SHA256SUMS} is published and \code{v1.0} is a signed tag.
  \item The design review is written and reviewed.
\end{checklist}
\end{milestonebox}

\begin{questionbox}
\begin{enumerate}
  \item Which parts of Kagan trust input from user programs, and how is each validated?
  \item Give an example of a TOCTOU bug that could exist in \code{execve}, and explain how copying first prevents it.
  \item Where would MINIX 3 put the Kagan ATA driver, and what would change about a driver bug's consequences?
  \item What is the most expensive operation measured, and where does the time go?
  \item Which of \MOS chapter 12's principles did Kagan follow best, and which did it violate most?
  \item What would break first if Kagan ran on two CPUs tomorrow?
\end{enumerate}
\end{questionbox}

\section{Summary}

Kagan 1.0 boots on a PC, protects processes from each other,
schedules them preemptively, manages virtual memory with copy-on-write, stores files in a real on-disk format, and runs
a shell with pipes. Every path through the system, from a key press to a character on the screen, from \code{fork}
to \code{exit}, from \code{write} to the disk, is implemented, tested and documented. Further work continues in the
engineering journal, in the extension tracks, and in Stage V.

\nextphase{16}{Portability: an Architecture Layer}
\booklegend
\end{document}
